\section{Extremal prime sums on a prescribed sparse sequence}\label{sec:sampling}

The earlier ProveIt note \emph{Torsion and prime sampling} established a
one-sided zero criterion on prescribed power grids and on much thinner
grids for a residual set of bases. It asked how much survives for a
prescribed base. Here one and the same explicit set of integer exponents
works for every real base $b>1$, including $b=2$. It recovers the exact
upper limit under critical-line RH and detects an attained rightmost
zero without any zero hypothesis. The set of exponents is
\[
 M_j=2^{2^{j^2}},\qquad
 \mathcal E=\bigcup_{j\ge1}\{M_j,2M_j,\ldots,jM_j\}.
\]
Among the cutoffs $x=b^e$, $e\in\mathcal E$, at most
$O_b(1+\log\log\log X)$ lie below $X$ as $X\to\infty$.
Indeed $b^{M_j}\le X$ implies $j^2=O_b(1+\log\log\log X)$,
and the first $j$ blocks contain at most $j(j+1)/2$ terms.

For a number field $K$, write $\Lambda_K(n)$ for the coefficients of
$-\zeta_K'/\zeta_K$ and $\gamma_K$ for its Euler--Kronecker constant,
with the convention $\zeta_K'/\zeta_K(s)=-(s-1)^{-1}+\gamma_K+O(s-1)$.
Take the ramp used in the preceding note, modeled on the \ghsource{grh/src/New/Reduction/OneSidedBound.lean}{prime remainder in the Lean proof},
\[
 R_K(x)=\sum_{n\le x}\Lambda_K(n)(1/n-1/x)-\log x+1+\gamma_K.
\]
For each nontrivial zero $\rho=\sigma+i\gamma$ put
$c_\rho=1/[\rho(1-\rho)]$ and let $\Theta$ be the supremum of the
real parts of all these zeros. All sums count multiplicity.

\paragraph{Theorem.}
If every nontrivial zero of $\zeta_K$ has real part $1/2$, then
for every $b>1$,
\[
 \lim_{j\to\infty}\max_{1\le k\le j}
 b^{kM_j/2}R_K(b^{kM_j})
 = \sum_{\rho}\frac{1}{1/4+(\operatorname{Im}\rho)^2}
 =: S_K.
\]
For $K=\mathbb Q$ the repository's RH theorem applies and
$S_{\mathbb Q}=2+\gamma_{\mathbb Q}-\log(4\pi)
=0.0461914\ldots$.
Unconditionally for every number field, the limsup of
$x^{1-\Theta}R_K(x)$ on $x=b^e$, $e\in\mathcal E$, is strictly positive
if and only if some zero has real part $\Theta$. When no zero attains
$\Theta$, that normalized expression tends to zero on the entire ray.

\paragraph{Proof.}
The smoothed explicit formula is
\[
 R_K(x)=\sum_\rho c_\rho x^{\rho-1}
       +O_K(x^{-1}\log(2x)),\qquad x\ge2 .
\]
For the normalization, the uncentered sum is $(1-1/x)\Phi_K(x)$
in Ihara's notation. His explicit formula (1.2.1), the archimedean
term (1.2.4), and the constant identity (1.3.11) give the displayed
formula. Equivalently, multiplying $-\zeta_K'/\zeta_K(s)$ by
$x^{s-1}/[s(s-1)]$ gives the residues $c_\rho x^{\rho-1}$
and $\log x-1-\gamma_K$ at the zeros and the pole at $1$,
respectively. The zero count $O_K(T\log(2T))$ implies
$\sum_\rho|c_\rho|<\infty$.
By dominated convergence, the sum over $\operatorname{Re}\rho<\Theta$
is $o(1)$ after multiplying by $x^{1-\Theta}$. The normalized
error also tends to zero. Thus the only surviving sum is
$\sum_{\operatorname{Re}\rho=\Theta}c_\rho
\exp(i\gamma\log x)$.

Fix any finite list of its $d$ ordinates, set
$\theta_i=\gamma_i\log b/(2\pi)$, and take an integer $Q\ge2$.
The fractional parts of $0,M_j\theta,\ldots,Q^dM_j\theta$ occupy
$Q^d$ half-open boxes of side $1/Q$ in $[0,1)^d$. The difference between
two points in one box supplies $1\le k\le Q^d$ with
$\|kM_j\theta_i\|_{\mathbb R/\mathbb Z}\le1/Q$ for all $i$.
For every $j\ge Q^d$ this $k$ belongs to the specified block.
In fact a positive fraction of each large block works. For $N=j+1$
points and $B=Q^d$ boxes, the number of same-box pairs is at least
$(N^2/B-N)/2$ by Cauchy--Schwarz. Each positive difference occurs
in at most $N$ pairs. Thus there are at least $(N/B-1)/2$ distinct
suitable $k$ in the block.
Absolute summability lets us first discard a tail of arbitrarily
small norm and then choose $Q$ to make all remaining phases arbitrarily
close to one. Hence a subsequence of the normalized values tends to
$\sum_{\operatorname{Re}\rho=\Theta}c_\rho$. This sum is real by
conjugation. Every summand has positive real part because
$\operatorname{Re}[\rho(1-\rho)]=\sigma(1-\sigma)+\gamma^2>0$.
This proves the assertion about attainment, including the empty case.
On the critical line every $c_\rho$ is positive real. The triangle
inequality supplies the matching upper bound on every block. The
value of the zero sum is explicit in field invariants. For signature
$(r_1,r_2)$, discriminant $d_K$ and Euler's constant $\gamma$,
Ihara (1.2.2), (1.2.3) and (1.3.11) give
\[
 S_K=2+2\gamma_K+\log|d_K|-r_1(\gamma+\log(4\pi))
                    -2r_2(\gamma+\log(2\pi)).
\]
Taking $K=\mathbb Q$ gives the stated numerical constant. Moreover,
for every fixed $\varepsilon>0$, a positive fraction of every sufficiently
large block comes within $\varepsilon$ of $S_K$ on RH. This follows by
fixing the tail and $Q$ for that tolerance in the same-box count.

\paragraph{Consequence and scope.}
For an arbitrary field $K$, a finite upper bound on
$\sqrt{x}R_K(x)$ on this grid rules out an attained rightmost zero
with real part greater than $1/2$. Indeed, on the subsequence just
constructed, the extra factor $x^{\Theta-1/2}$ tends to infinity.
The case where off-line real parts approach their supremum strictly
from below remains open for this test. The earlier question also specified
a different, preassigned sequence of exponents $\exp(n/\log(n+2))$.
The block recurrence proved here addresses the fixed-base issue on
the displayed grid. Sparsity counts cutoffs and conveys no
computational savings in evaluating a prime sum at such a large cutoff.

\smallskip\noindent
For the explicit formula see Y.~Ihara, \href{https://www.kurims.kyoto-u.ac.jp/~kenkyubu/emeritus/ihara/Publications-and-Recent-Preprints/RecentArticles/pdf-files/1A.pdf}{\emph{On the Euler--Kronecker constants of global fields and primes with small norms}},
Part~1, \S\S1.2--1.3.
The antecedent ProveIt argument and its open question appear in
\href{https://github.com/VladimirReshetnikov/ProveIt/blob/d31039a48d963692e5a5176f47b6603fc9e9d9cf/docs/incoming/grh-torsion-and-prime-sampling.zip}{\texttt{grh-torsion-and-prime-sampling.zip}},
in the section on sparse scales.
